% Version 0.27. Actual symplectic local and cohomological installation.
% Pinned source: 7cbf7a0c32fca1badf05931879a37c3f8a2b03b4.
\section{Installing the arithmetic symplectic deformation data}
\label{sec:symplectic-local-installation}

\status{DX1--DX8 close the three actual arithmetic inputs to
\texttt{exModel}: \texttt{PhiCond}, \texttt{Nums}, and an original
\texttt{LocBase}.  The ordinary and finite-inertia charts are proved
by continuous cohomology and explicit matrix deformations.
Their smoothness is at the characteristic-zero origin, precisely
the source's \texttt{SmoothPt}; integral formal smoothness is not
asserted.  The full \texttt{BlochKatoLocal} package is not assumed
or promoted.  The original \texttt{Good} datum remains the input,
supplied by the separate arithmetic construction.}

\subsection{Original fields, characters and roots}

Fix the source's original algebraic closure $\Omega/\mathbb Q$,
odd prime $\ell$, coefficient DVR $O=c.O$, fraction field $E$,
and finite residue field $k$.  Fix a \texttt{Raw} datum $D$
satisfying \texttt{Good}, without replacing its places or
decomposition subgroups.  Write
\[
 \begin{aligned}
 \chi&=\chi_\ell,&
 \alpha&=\chi\omega^{-1}\xi,&
 \Psi&=\chi^e\psi,\\
 \beta&=\Psi\alpha^{-1},&
 \theta&=\alpha\Psi,&
 r_0&=\operatorname{diag}(\theta,\Psi,\alpha,1).
 \end{aligned}
 \tag{DX.1}
\]
Here $e\ge2$, $\psi,\omega,\xi$ have finite image, $\xi$ is
unramified at every member of $T=\{\ell\}\cup S_b$, and
$\omega$ is the specified residual cyclotomic lift.
All further arithmetic facts used below are the original
\texttt{facts(D,Good)} supplied by AX6.

Let $G_v=D.Gp(v)$ and $I_v=D.Ip(v)$, and retain the original
chosen $s_0\in I_\ell\cap\ker\psi\cap\ker\omega$.
Its cyclotomic value has infinite order.  Thus the four eigenvalues
of $r_0(s_0)$ are distinct, since their exponents are
$e+1,e,1,0$.  At every $v\in T$ the original $\sigma_v=D.sg(v)$
has the stated residual separation; at $\ell$,
$\chi(\sigma_\ell)=1$.

Use the weights
\[
 \begin{aligned}
 a&=(1,0),& b&=(-1,1),& p&=(0,1),& t&=(1,1),\\
 w_0&=(1,1),& w_1&=(0,1),& w_2&=(1,0),& w_3&=0 .
 \end{aligned}
 \tag{DX.2}
\]
The letter $p$ in this weight list denotes $w_{\rm Ps}$;
prime subscripts below are written $q$ when ambiguity is possible.
With the original alternating form
\[
 J=\begin{pmatrix}
 0&0&0&1\\0&0&1&0\\0&-1&0&0\\-1&0&0&0
 \end{pmatrix},
\]
the Lie algebra $\mathfrak{sp}_4$ has the following weight basis:
\[
\begin{array}{c|c|c}
\text{weight}&\text{basis matrix}&\text{adjoint character}\\ \hline
0&\operatorname{diag}(1,0,0,-1),\
       \operatorname{diag}(0,1,-1,0)&1\\
a&E_{01}-E_{23}&\alpha\\
p&E_{02}+E_{13}&\Psi\\
b&E_{12}&\beta\\
t&E_{03}&\theta\\
-a&E_{10}-E_{32}&\alpha^{-1}\\
-p&E_{20}+E_{31}&\Psi^{-1}\\
-b&E_{21}&\beta^{-1}\\
-t&E_{30}&\theta^{-1}
\end{array}
 \tag{DX.3}
\]
Direct multiplication verifies $Y^tJ+JY=0$ for these matrices;
that equation also determines all other entries from the displayed
ten.  The table is therefore valid over $O$, $E$, and $k$.
Put
\[
 d_P(j)=
 \begin{cases}2&j=0,\\1&j\in\{\pm a,\pm b,\pm p,\pm t\},\\0&\text{otherwise},
 \end{cases}
 \quad
 d_L(j)=
 \begin{cases}2&j\in\{0,a,b,p,t\},\\1&j\in\{-a,-b,-p,-t\},\\0&\text{otherwise}.
 \end{cases}
 \tag{DX.4}
\]
These are literally \texttt{dimP} and \texttt{dimL}.

\subsection{The exact finite and rational cohomology inputs}

% BEGIN DX1
\begin{proposition}[DX1: original local duality, global Euler formula and normalized finite duality]
\label{claim:DX1}
The native \texttt{LocalFiniteness}, \texttt{LocalTateDuality},
\texttt{LocalEulerChar}, \texttt{GlobalFiniteness} and
\texttt{GlobalEulerChar} interfaces follow from the finite local
and global theorems stated below, with their original coefficient
modules, places and cup products.  LD3 supplies the literal
\texttt{PoitouTateNormalised} by forgetting its additional
conjugation-compatibility field.
These universally quantified interface assertions range over
every prime, number field and finite coefficient module
allowed by their definitions; they do not inherit an
odd-prime restriction from the later good-datum application.

In particular, for a finite $k[G_{\mathbb Q,S}]$-module $M$,
with $\ell\in S$, write $h^i=\dim_kH^i$ and let $c$ be the
original complex conjugation.  Then
\[
 h^1_{\rm glob}+ \dim_k M^{c=1}
       =h^0_{\rm glob}+h^2_{\rm glob}+\dim_k M,
 \tag{DX.5}
\]
and on the original $G_v$,
\[
 h^1_v=h^0_v+h^2_v+
                  {\bf1}_{v=\ell}\dim_kM .
 \tag{DX.6}
\]
For a continuous unit-valued character
$\kappa:G_\ell\to O^\times$, rational continuous cohomology is
finite dimensional and
\[
 \dim_EH^2(G_\ell,E(\kappa))
       =\dim_E H^0(G_\ell,E(\kappa^{-1}\chi)),\qquad
 \dim_EH^1=\dim_EH^0+\dim_EH^2+1 .
 \tag{DX.7}
\]
All coefficient twists in these formulas use the original
cyclotomic action.
\end{proposition}

\begin{proof}
The named classical inputs are finite local Tate duality and the
local Euler--Poincare formula, global finiteness, and the global
Euler--Poincare formula.  Precise primary references are given
at the end of the section.  We spell out their translation,
including the rational passage needed for the charts.

For a valuation subring $A$ over a finite place $v$ of an arbitrary
number field $F$, the original decomposition group is the absolute
Galois group of the corresponding henselization.  Finite separable
extensions of that henselian valued field and of its completion
correspond, preserving the embeddings into the chosen separable
closures.  This gives the topological Galois-group identification
with $G_{F_v}$ used in KB6.  Its inertia subgroup is the original
inertia subgroup.  It also preserves roots of unity: the action
on $\mu_{\ell^r}$ is exactly the pullback of
$\operatorname{cycMod}(F,\Omega,\ell,r)$.
Precomposition consequently gives isomorphisms of the actual
continuous cochain complexes, commuting with restriction and cup
products.  No unrelated coefficient or local field is substituted.

For $R=\mathbb Z/\ell^r$ and any finite, possibly nonfree, $R$-module
$M$, evaluation and a chosen primitive root identify the source's
twisted dual with $\operatorname{Hom}(M,\mu_{\ell^r})$.
The finite local theorem gives finiteness, vanishing above degree
two, and perfect pairings in degrees $(0,2),(1,1),(2,0)$.
Kummer theory identifies $H^2(R(1))$ with the $\ell^r$-torsion
in the local Brauer group.  Its invariant identifies that target
with $R$.  Thus perfection gives the source's maps into the
dual \emph{before} selecting the invariant as well.  The reversed
degree-one map is also bijective: graded commutativity changes
only a sign.  This proves every native local-duality clause.
It is the same finite-coefficient argument as LK1--LK3,
now transported to the actual decomposition groups.

The finite local Euler formula is
\[
 \frac{|H^0(F_v,M)|\,|H^2(F_v,M)|}{|H^1(F_v,M)|}
                   =|\mathcal O_{F_v}/|M|\mathcal O_{F_v}|^{-1}.
 \tag{DX.8}
\]
If $v\nmid\ell$ the right side is one.
If $v\mid\ell$ it is $|M|^{-[F_v:\mathbb Q_\ell]}$:
compute successively along the residue-field and ramification
filtrations, exactly as in LK4.  This proves the native
\texttt{locDeg} exponent even when $M$ is not free.

For the global assertions, the quotient by the closed normal
subgroup generated by the omitted inertia groups is precisely
$G_{F,S}$: the assertion is tested on finite quotients, where it
is the definition of being unramified outside $S$.
The classical finite set includes all archimedean places; this
places no additional restriction on that quotient.
Because $S$ contains the places above $\ell$, $|M|$ is a unit
in the ring of $S$-integers.
The global finiteness theorem therefore applies.  Its
high-degree comparison is with the real-place groups, which
vanish for odd $\ell$, proving the exact extra clause in
\texttt{GlobalFiniteness}.

The global Euler formula reads
\[
 |H^0|\,|H^2|\,|M|^{[F:\mathbb Q]}
   =|H^1|
     \prod_{v\ {\rm real}}|M^{c_v=1}|\,
                       |M|^{r_2(F)}.
 \tag{DX.9}
\]
Indeed a real normalized absolute value sends $|M|$ to $|M|$,
and a complex one sends it to $|M|^2$.
This is exactly \texttt{archH0Card}; it also covers $\ell=2$
in the stated Euler interface, without claiming high-degree
vanishing at real places in that case.
Milne explicitly defines the Euler expression immediately
before I.5.1 by the truncated ratio $|H^0||H^2|/|H^1|$,
including the real-place, $2$-primary case where higher
cohomology need not vanish.

The source's \texttt{IsLocalInvariants} contains bijectivity and
the uniformizer normalization required by
\texttt{PoitouTateNormalised}, together with conjugation
compatibility.  The rest of their quantified conclusions is
identical: the invariant family is chosen before $S$, the places,
the descended cyclotomic character and $M$.
Forget that one extra field in LD3.
This is a literal implication between the two definitions,
not an identification based on their names.
For a finite coefficient field $k$, restricting scalars to
$\mathbb F_\ell$ and using its nondegenerate trace identifies
the ordinary finite dual with the $k$-linear twisted dual.
All dimensions over $\mathbb F_\ell$ acquire the same factor
$[k:\mathbb F_\ell]$.  Taking cardinalities in (DX.8)--(DX.9)
gives (DX.5)--(DX.6), while the same trace identification
preserves the finite Poitou--Tate sequence.

For completeness, let $T=O(\kappa)$ and choose a uniformizer
$\varpi$ of $O$.  Continuous cochains with values in $T$ are
the inverse limit of those with values in $T/\varpi^nT$.
The cochain transition maps are onto: a locally constant map
into a finite quotient has a locally constant lift.
All finite-level cohomology groups are finite.  Their inverse
systems satisfy the Mittag--Leffler condition, so
\[
 H^i(G_\ell,T)=\varprojlim_nH^i(G_\ell,T/\varpi^nT).
 \tag{DX.10}
\]
The connecting sequence embeds
$H^i(T)/\varpi H^i(T)$ into the finite group $H^i(T/\varpi T)$.
Choose finitely many lifts of generators of this quotient.
Successive $\varpi$-adic approximation shows that they generate
$H^i(T)$: the corrections tend to zero in every finite-level
projection, and the image of a finite power of compact $O$
is closed.  Thus these compact cohomology groups are finite
$O$-modules.

A continuous cochain into $E(\kappa)$ has compact, hence bounded,
image on the compact source $G_\ell^i$; one power of $\varpi$
puts it into $T$.  Consequently rational cohomology is
$H^i(T)\otimes_OE$.
The compatible finite Tate pairings identify the Pontryagin dual
of $H^2(T)$ with the invariants of its discrete dual lattice.
Using the $\mathbb Z_\ell$-dual lattice rather than an
unjustified self-duality of $O$ introduces the codifferent,
but after tensoring with $E$ its character is still
$\kappa^{-1}\chi$.
If this character is not one, one nonzero scalar
$\kappa(g)^{-1}\chi(g)-1$ kills all its discrete invariants;
they are therefore finite.  If it is one, the invariants are
the entire divisible rank-one module.  The $E$-dimension of
$H^2$ is accordingly zero or one, giving the first formula
in (DX.7).

Finally the exact sequences
\[
 0\to H^i(T)/\varpi^nH^i(T)
 \to H^i(T/\varpi^nT)
 \to H^{i+1}(T)[\varpi^n]\to0
\]
show that the logarithm to base $|k|$ of the middle cardinality
is $n\operatorname{rank}_OH^i(T)+O(1)$.
Apply (DX.8), divide by $n$, and let $n$ increase.
The coefficient length of $T/\varpi^nT$ is $n$, regardless of
the ramification of $O/\mathbb Z_\ell$.
This proves the second formula in (DX.7) with the displayed
coefficient-field normalization.
\end{proof}
% END DX1

% BEGIN DX2
\begin{proposition}[DX2: actual global and local \texttt{PhiCond}]
\label{claim:DX2}
For the original good datum,
$\Phi_\ell(D.Gam)$ holds.  It also holds for every original
decomposition group $D.Gp(v)$ at a finite prime $v$.
These conclusions require no Bloch--Kato theorem.
\end{proposition}

\begin{proof}
We use a concrete finite-cochain argument that avoids changing the
original group.  Let $\Gamma$ be a profinite group for which
$H^1(\Gamma,M)$ is finite for every finite continuous
$\mathbb F_\ell[\Gamma]$-module.  If $H\le\Gamma$ is open,
the coset set $X=\Gamma/H$ is finite and discrete.
Choose representatives $s(x)$, taking $s(H)=1$.
Give $P=\mathbb F_\ell^X$ the permutation action
$(g f)(x)=f(g^{-1}x)$.
For a continuous additive homomorphism $u:H\to\mathbb F_\ell$,
put
\[
       c_u(g)(x)=u\bigl(s(x)^{-1}g\,s(g^{-1}x)\bigr).
 \tag{DX.11}
\]
The argument belongs to $H$.  Multiplication of two such
arguments proves $c_u(gh)=c_u(g)+g c_u(h)$.
There are only finitely many cosets, so this is a continuous
cocycle.  Evaluation at $(h,H)$ gives $u(h)$, making
$u\mapsto c_u$ injective.
The cocycle group $Z^1(\Gamma,P)$ is finite: its quotient
$H^1(\Gamma,P)$ is finite and its coboundary subgroup is an
image of finite $P$.  Hence there are only finitely many $u$.
Open-kernel homomorphisms to the finite target are continuous,
so this is exactly the source's \texttt{PhiCond}.

For $D.Gam=G_{\mathbb Q,S}$, the finite set $S$ consists of the
original rational primes and contains $\ell$.  Apply global
finiteness from DX1.  For $D.Gp(v)$ use the original
decomposition/completion identification in DX1 and finite
local cohomology, also for modules of order prime to $v$.
This proves both assertions with the stated original topologies.
\end{proof}
% END DX2

\subsection{The complete framed numerical calculation}

% BEGIN DX3
\begin{proposition}[DX3: finite Selmer duality and the actual framed Euler identity]
\label{claim:DX3}
Finite Poitou--Tate duality, with DX1's matched invariants, gives
the literal \texttt{GreenbergWiles} identity for arbitrary finite
$\mathbb Z/\ell^r$-modules and arbitrary local submodules, when
$\ell$ is odd.

For the original rational field and a frame with
$S=T\sqcup Q$, $\ell\in T$, distinguished $v_0\in T$, put
$N_q=I_q\cap K$.  Suppose $K$ is normal, the action on $N_q$
is trivial, and $[I_q:N_q]$ is prime to $\ell$.
For a finite $k$-representation $M$ let $H^1_F(M)$ and
$\operatorname{Ob}_F(M)$ be the exact cochain spaces defined in
FP.4--FP.6, with the literal intersection quotient for
$\operatorname{Ob}$.  They are finite dimensional and
\[
 \dim\operatorname{Ob}_F(M)-\dim H^1_F(M)
       =\dim M^{c=1}+\dim H^0(G_{\mathbb Q,S},M^\vee(1)).
 \tag{DX.12}
\]
\end{proposition}

\begin{proof}
First fix an arbitrary number field and odd $\ell$, a finite
coefficient module $A$, and local submodules
$L_v\subset H^1(F_v,A)$.  Let $L_v^\perp$ be their exact
annihilators under the original local cup product.
Finite perfect duality and the middle and final terms of
LD3's nine-term sequence give the exact sequence
\[
\begin{aligned}
0\to\operatorname{Sel}_{L}(A)&\to H^1(G_{F,S},A)
 \to\bigoplus_{v\in S}H^1(F_v,A)/L_v\\
 &\to\operatorname{Sel}_{L^\perp}(A^\vee(1))^\vee
 \to H^2(G_{F,S},A)\\
 &\to\bigoplus_{v\in S}H^2(F_v,A)
 \to H^0(G_{F,S},A^\vee(1))^\vee\to0 .
\end{aligned}
 \tag{DX.13}
\]
To see the only modified middle term, divide the local
degree-one term in the Poitou--Tate sequence by $\bigoplus L_v$.
Its annihilator in the global dual cohomology is exactly the
inverse image of $\bigoplus L_v^\perp$.  Restriction of
functionals identifies the quotient of its dual by the
annihilator with the displayed Selmer dual.
Dualizing finite modules here is exact: the
$\mathbb Z/\ell^r$-linear dual is the group-character dual
under $a\mapsto a/\ell^r$, including nonfree modules.

Take cardinalities in (DX.13).  Apply (DX.8)--(DX.9) and
$\sum_{v\mid\ell}[F_v:\mathbb Q_\ell]=[F:\mathbb Q]$.
All powers of $|A|$ cancel, leaving
\[
\begin{aligned}
|\operatorname{Sel}_L(A)|\,|H^0(A^\vee(1))|
 \prod_{v\in S}|H^0(F_v,A)|\,\operatorname{archH0Card}(A)
={}&|\operatorname{Sel}_{L^\perp}(A^\vee(1))|\\
&{}\cdot |H^0(A)|\prod_{v\in S}|L_v|.
\end{aligned}
 \tag{DX.14}
\]
This is the full native \texttt{GreenbergWiles} formula.
All local maps and annihilators were taken on the original
places; no alternate normalization is required.

We now prove the framed identity without treating the framed
obstruction space as an unnamed Selmer complex.
Write $U_q=H^1_{\rm ur}(G_q,M;N_q)$, and define
\[
 \begin{aligned}
 \operatorname{Sel}&=\ker\left[
 H^1(G_{\mathbb Q,S},M)\to
  \bigoplus_{v\in T}H^1(G_v,M)
       \oplus\bigoplus_{q\in Q}H^1(G_q,M)/U_q\right],\\
 \operatorname{Sha}^2&=\ker\left[
 H^2(G_{\mathbb Q,S},M)\to
                     \bigoplus_{v\in S}H^2(G_v,M)\right].
 \end{aligned}
 \tag{DX.15}
\]
The map from $H^1_F$ to $\operatorname{Sel}$ sends a cocycle
and its local primitives to its global cohomology class.
It is onto: choose local primitives at every $v\in T$ and
subtract the coboundary of the primitive at $v_0$.
At $q$, vanishing of the restricted class on $N_q$ is actual
vanishing of that cocycle, since $N_q$ acts trivially.
The kernel is parametrized by the local invariant vectors,
modulo one common global invariant vector.  Explicitly, for
a global coboundary $d^0m$, its primitives have the form
$m+b_v$, with $b_v\in M^{G_v}$; the normalization at $v_0$
determines $m=-b_{v_0}$, and the remaining redundancy is
$M^{G_{\mathbb Q,S}}$.  Thus
\[
 \dim H^1_F+h^0_{\rm glob}
             =\dim\operatorname{Sel}+\sum_{v\in T}h^0_v .
 \tag{DX.16}
\]

The map $\operatorname{Ob}_F\to H^2(G_{\mathbb Q,S},M)$
sends a triple to its first cocycle class.
Its image is exactly $\operatorname{Sha}^2$, since the local
components are precisely chosen primitives for its restrictions.
If that first class is zero, subtract its global primitive.
The remaining local entries are degree-one cocycles.
Quotienting by FP.6 divides the $T$-components by coboundaries,
the $Q$-components by $U_q$, and then divides by the image
of global $H^1$.  Conversely every such boundary gives
exactly those identifications.  Hence
\[
0\to\operatorname{coker}\left[
 H^1_{\rm glob}\to
 \bigoplus_T H^1_v\oplus\bigoplus_Q H^1_q/U_q\right]
 \to\operatorname{Ob}_F\to\operatorname{Sha}^2\to0 .
 \tag{DX.17}
\]
This proves finiteness as well as the dimension formula
\[
 \dim\operatorname{Ob}_F+h^1_{\rm glob}+\sum_Q\dim U_q
 =\dim\operatorname{Sha}^2+\sum_T h^1_v+\sum_Q h^1_q
                              +\dim\operatorname{Sel}.
 \tag{DX.18}
\]

Restriction from $I_q$ to $N_q$ has zero kernel on $H^1$:
inflation--restriction identifies that kernel with
$H^1(I_q/N_q,M)$, which is zero by averaging over the
finite group of order prime to $\ell$.
Thus $U_q$ is the usual unramified subgroup.
The quotient $G_q/I_q=\widehat{\mathbb Z}$ gives
$U_q=M^{I_q}/(\operatorname{Frob}_q-1)M^{I_q}$, so
$\dim U_q=h^0_q$ by elementary rank-nullity.
Since $q\ne\ell$, (DX.6) gives $h^1_q-\dim U_q=h^2_q$.
The sum of (DX.6) over $T$ contributes exactly one
$\dim M$, because $\ell\in T$.

Finally the last three nonzero terms of finite Poitou--Tate give
\[
 \dim\operatorname{Sha}^2+\sum_{v\in S}h^2_v
       =h^2_{\rm glob}+\dim H^0(M^\vee(1)).
 \tag{DX.19}
\]
Subtract (DX.16) from (DX.18), then use these local identities,
(DX.19) and (DX.5).  The result is exactly (DX.12).
\end{proof}
% END DX3

% BEGIN DX4
\begin{proposition}[DX4: the literal original \texttt{Nums}]
\label{claim:DX4}
Every original good $D$ satisfies $D.\texttt{Nums}$.
More explicitly, for every descent $\bar r_S$ of the specified
$\bar r$ to $D.Gam$, every supplied normalization proof, and
every $j\in\mathbb Z^2$, the actual framed obstruction space
is finite dimensional and
\[
 \begin{aligned}
 \dim_k\operatorname{Ob}_F(\operatorname{ad}_j)
   &=\dim_kH^1_F(\operatorname{ad}_j)+
      \begin{cases}2&j=0,\\1&j=p\text{ or }j=-p,\\0&\text{otherwise};
      \end{cases}\\
 H^1_F(\operatorname{ad}_j)&=0
                         \quad(j\notin\{0,p,-p\}).
 \end{aligned}
 \tag{DX.20}
\]
The frame, all inertia conditions, coefficient field and complex
conjugation are the original ones.
\end{proposition}

\begin{proof}
Let $K=\ker\bar r$.  AX6 makes it open and normal.
Its finite quotient embeds in $(k^\times)^4$, so its order
and every inertia index $[I_q:I_q\cap K]$ are prime to $\ell$.
The frame is exactly $T=\{\ell\}\cup S_b$, $Q=S\setminus T$,
with $N_q=I_q\cap K$ and $v_0=\ell$.
The set $T\sqcup Q$ is the original $S$.
Thus all hypotheses in DX3 hold, including the exact original
local triviality at $N_q$.

For the weight-$j$ module in (DX.3), the action is scalar by
\[
       \gamma_j=\overline{\alpha^{j_1}\Psi^{j_2}}.
 \tag{DX.21}
\]
Indeed an $(u,v)$ entry transforms by
$\overline{\operatorname{ps}_u/\operatorname{ps}_v}$.
The displayed basis therefore proves both the action and
dimension $d_P(j)$, independently of a choice of descent:
the quotient map to $D.Gam$ is onto.

The twisted dual has no global invariant unless its scalar
character is cyclotomic.  For $j=0,\pm p$ this cannot happen:
the original $c=D.cc$ acts by $1$ on $\gamma_j$ and by $-1$
on the residual cyclotomic character.
For $j\in\{\pm a,\pm b,\pm t\}$, use the original
$\sigma_\ell$.  The cyclotomic value is $1$, whereas
AX6's residual separation makes $\gamma_j(\sigma_\ell)\ne1$.
For all other weights the module is zero.
Consequently the last term of (DX.12) is always zero.

At complex conjugation,
$\gamma_j(c)=(-1)^{j_1}$.
Thus its fixed-space dimension is two at zero, one at
$\pm p$, and zero at the six remaining roots.
Formula (DX.12) proves the equality and finiteness in (DX.20).

It remains to prove the asserted vanishing of the actual framed
cocycles.  At a weight outside (DX.3) the coefficient space
is zero.  At an odd root it is one dimensional and the
appropriate \texttt{Good.xvan} applies to
$\alpha,\beta,\theta$ or its inverse.
Take a framed pair $(z,(y_v))$ and pull $z$ back to the
original $G_{\mathbb Q}$.
It vanishes on inertia outside $S$, since it factors through
$G_{\mathbb Q,S}$, and it vanishes on $I_q\cap K$ at
the original chosen place for every $q\in Q$.
The same holds at every conjugate place:
\[
 z(hnh^{-1})=\gamma_j(h)z(n)
                           +(1-\gamma_j(n))z(h)=0
       \quad(n\in K).
 \tag{DX.22}
\]
This is exactly the original \texttt{UnrT} condition,
with no loss of valuations.
The framed primitives give the required coboundaries at $T$.
By \texttt{xvan}, $z(g)=(\gamma_j(g)-1)m$ globally.
At $v_0=\ell$ the original primitive is zero.
Evaluate at $\sigma_\ell$, where $\gamma_j\ne1$, to obtain
$m=0$.  Then $z=0$.
At every other $v\in T$, evaluation at the original
$\sigma_v$ and its residual separation gives $y_v=0$.
Hence the framed space itself is zero, not only its
cohomology quotient.  This verifies every quantifier and
field of $D.\texttt{Nums}$.
\end{proof}
% END DX4

\subsection{From integral universal representations to actual weighted charts}

% BEGIN DX5
\begin{proposition}[DX5: the integral ring and the characteristic-zero chart]
\label{claim:DX5}
Let $\Gamma$ be an original decomposition group from DX2.
Let $\mathcal P$ be either $\operatorname{eqL}$ at $\ell$ or
$\operatorname{eqP}$ at $q\in S_b$.
There is an integral universal ring $R\in\mathcal C_O$ for the
actual continuous lifts satisfying $\mathcal P$, with a weighted
presentation $W$, all variable weights in
$\{0,\pm a,\pm b,\pm p,\pm t\}$, and an origin
$x:R\to O$ giving exactly $r_0|_\Gamma$.
Every universal $(i,j)$ matrix entry has weight $w_i-w_j$.

Suppose the functor of continuous $\mathcal P$-lifts of
$r_0|_\Gamma$ over augmented finite local $E$-algebras is
unobstructed for small extensions, and its weight-$j$ tangent
space has dimension $n(j)$.
Then the original $W.\!A$, with its actual origin and grading,
satisfies
\[
                \operatorname{SmoothPt}(W.\mathrm{gr},x,n).
 \tag{DX.23}
\]
More precisely, the resulting map
$\phi:R\to E[[t_1,\ldots,t_d]]$ has the source's
$\operatorname{IsCompl}(\ker\phi,x,\phi)$ property, where
$d=\sum_j n(j)$; some $y\in R$ with $x(y)\ne0$
annihilates $\ker\phi$.
No integral smoothness assertion is needed.
\end{proposition}

\begin{proof}
By DX2, FP2 constructs the universal continuous matrix lift.
The prescribed conditions are actual equations in its entries.
For $\operatorname{eqP}$ they are
\[
 \rho(g)^tJ\rho(g)=\theta(g)J,\qquad
 \rho(n)=1\quad(n\in I_q\cap\ker\psi).
 \tag{DX.24}
\]
For $\operatorname{eqL}$ add the specified characteristic
polynomial of $\rho(s_0)$ and the source's equations
\[
 \begin{aligned}
 P_j(\rho(s_0))\,\rho(g)\,P_i(\rho(s_0))&=0
                                      &&(i<j),\\
 P_i(\rho(s_0))(\rho(s)-\operatorname{ps}_i(s))P_i(\rho(s_0))
                                      &=0 &&(s\in I_\ell),
 \end{aligned}
 \tag{DX.25}
\]
where
$P_i(X)=\prod_{h\ne i}(X-\operatorname{ps}_h(s_0))$.
These are the unnormalized products in \texttt{Ord};
no eigenvalue difference has been inverted over $O$.
Quotient the universal ring by the ideal generated by all these
entries.  That ideal is closed and finitely generated by
Noetherianity, even though the list of equations was infinite.
The quotient is proper, because the original diagonal $r_0$
satisfies every equation.  Universality of the quotient follows
by applying the equations entry by entry under an arbitrary
$O$-algebra map.  Thus it is the actual integral
$\operatorname{Represents}(\mathcal P,\rho)$, for every
$\mathcal C_O$ target.

The two-dimensional diagonal torus acts by
$\rho(g)_{ij}\mapsto u^{w_i-w_j}\rho(g)_{ij}$.
The fixed multiplier is preserved, since
$w_i+w_{3-i}=(1,1)$.
Conjugation preserves the characteristic polynomial and all
products in (DX.25), as well as (DX.24).
Take finitely many universal entry differences
$\rho(g)_{ij}-(r_0(g))_{ij}$ generating the relative
cotangent space.  They generate the complete ring over $O$
by FP1.  The formal-torus argument in FP3 makes the kernel
of this presentation homogeneous and places every actual
matrix entry in its stated weight part.  This applies in
arbitrary residue characteristic; independence of formal
characters, rather than evaluation on the finite group
$k^\times$, is the relevant argument.
All selected differences vanish under evaluation at $r_0$.
Consequently the presentation has the actual origin $x$.
Every matrix-entry weight belongs to the displayed nine-element
set, including weight zero.

Put $\mathfrak a=\ker x$ and choose $\varpi\in O$.
The characteristic-zero completion at this origin is
\[
 \mathcal R=\varprojlim_m(R/\mathfrak a^m)\otimes_OE
        =\widehat{R[1/\varpi]}_{\mathfrak a R[1/\varpi]}.
 \tag{DX.26}
\]
The quotient $R/\mathfrak a^m$ is a finite $O$-module:
in the chosen presentation it is generated by monomials of
degree less than $m$ in the variables.
Thus $\mathcal R$ is a complete Noetherian local
$E$-algebra with residue field $E$.
Equivalently it is the quotient of $E[[z_1,\ldots,z_s]]$
by the extended presentation ideal.
These assertions are the Noetherian completion theorem,
applied to the actual presentation; in particular, completion
does not change any quotient by a power of the augmentation
ideal.

We verify the represented functor, since replacing integral
deformations by arbitrary characteristic-zero representations
without this step would be insufficient.
Let $B=E\oplus\mathfrak n_B$ be a finite local augmented
$E$-algebra, with $\mathfrak n_B^m=0$.
An $E$-algebra map $\mathcal R\to B$ produces the continuous
representation obtained from the integral universal one.
Indeed the variables land in $\mathfrak n_B$; the image of
$\mathfrak m_R^N$ lies in $\varpi^{N-m+1}$ times a fixed
finite $O$-module for $N\ge m-1$.  Hence the induced map is
continuous from the adic topology of $R$ to the
$\ell$-adic topology of finite-dimensional $B$.

Conversely let $\rho_B$ be a continuous lift of $r_0$ to $B$.
The differences of its entries from those of $r_0$ lie in
$\mathfrak n_B$.  They have compact image, because
$\Gamma$ is compact, and hence lie in some finite
$O$-lattice $L\subset\mathfrak n_B$.
The $O$-subalgebra $B_0=O[L]$ is finite over $O$:
products of $m$ elements of $L$ vanish, so only the finitely
many monomials of degree less than $m$ are required.
It is local with residue field $k$ and complete; its maximal
ideal is generated by $\varpi$ and its nilpotent augmentation
ideal.  Its induced topology from $B$ equals its adic topology.
Thus $\rho_B$ is an actual lift over $B_0\in\mathcal C_O$.
The polynomial equations hold in $B_0$ because $B_0\hookrightarrow B$
is injective.  Integral universality gives $R\to B_0\to B$.
The augmentation is $x$, by uniqueness on the matrix entries;
therefore this map factors uniquely through $\mathcal R$.
This argument uses only bounded subsets of the given
finite-dimensional $E$-algebra.  It makes no assertion that
a compact subset of an algebraic closure lies in a finite
extension.

Choose homogeneous elements $f_1,\ldots,f_d\in\mathfrak a$
whose images form an $E$-basis of
$\mathfrak m_{\mathcal R}/\mathfrak m_{\mathcal R}^2$.
They may be chosen in the original $R$: its homogeneous
variable images span this vector space, and denominators
in finite $E$-linear combinations can be cleared in $O$.
The dimension of each weight part is $n(j)$.
Here the weight agrees with the source's tangent convention:
a functional supported on weight $j$ gives precisely
an entry deformation in $\mathfrak{sp}_4(E)_j$.
No change from $j$ to $-j$ is made.

Successive approximation gives a surjection
\[
       E[[t_1,\ldots,t_d]]\longrightarrow\mathcal R,
                         \qquad t_i\longmapsto f_i .
 \tag{DX.27}
\]
Its kernel $J_0$ has no constant or linear term.
We prove it is zero, rather than inferring smoothness from
the dimension count.  If $J_0\ne0$, let $n\ge2$ be the
least order of a nonzero member and choose one of order $n$.
Writing $P=E[[t_1,\ldots,t_d]]$, the map
\[
   P/\mathfrak m_P^{n+1}\longrightarrow
                         P/(J_0+\mathfrak m_P^{n+1})
 \tag{DX.28}
\]
is a small extension: its kernel is contained in degree $n$
and is killed by the maximal ideal.
Unobstructedness would lift the natural $\mathcal R$-point
of the right side to the left side.
Under this lift each $t_i$ changes only in degree $n$.
Substitution therefore leaves the nonzero degree-$n$
initial form of the chosen relation unchanged, a
contradiction.  Hence (DX.27) is an isomorphism.
The chosen elements are homogeneous, so its inverse respects
the formal torus action and every weight part.
Composing $R\to\mathcal R$ with this inverse gives $\phi$.

Its constant coefficient is $x$.  Each $f_i$ maps to $t_i$,
so it gives the required separate nonzero linear coefficient.
For the injectivity clause, if $x(f)=0$ and $\phi(f)$ has
zero linear terms, the completion isomorphism modulo the
square shows that $f/1\in\mathfrak a^2R[1/\varpi]$.
Thus $cf\in\mathfrak a^2$ for some $0\ne c\in O$.
We must prove this in the source's graded subalgebra $W.\!A$,
not just in $R$.
For homogeneous $f$ of weight $j$, express
$cf=\sum_{u,v}a_{uv}z_uz_v$, where the $z_u$ are the
homogeneous variable images.
Project each $a_{uv}$ to weight $j-w_u-w_v$.
Homogeneity of the presentation makes this projection
well-defined in $R$, and leaves the equality unchanged.
All the resulting factors belong to $W.\!A$, and every
$z_u$ vanishes at $x$.
Therefore $cf$ belongs to
$\ker(x|_{W.A})^2$.
An arbitrary element of $W.\!A$ is a finite sum of homogeneous
parts.  Its vanishing constant and linear coefficients hold
part by part; multiplying the finitely many denominators
proves the same assertion for that element.
This is exactly membership in $\operatorname{satSq}(x|_{W.A})$.
The grading, constant, separate-linear-coefficient and
injectivity clauses of \texttt{IsChart} are now all verified.
Their precise parameter counts give (DX.23).

The stronger completion statement follows from the same
argument modulo $\mathfrak a^N$ for each $N$.
A truncated power series has finitely many $E$-coefficients;
clearing their denominators gives the required density.
If $\phi(f)$ lies in the $N$th power of the parameter ideal,
then $cf\in\mathfrak a^N$ for some $0\ne c\in O$;
this is stronger than the stated
$cf\in\mathfrak a^N+\ker\phi$.
Finally localize $R$ at the prime $\mathfrak a$.
The completion in (DX.26) is also its maximal-ideal completion.
Krull intersection makes $R_{\mathfrak a}$ inject into that
completion.  Hence each $f\in\ker\phi$ is killed by some
$s\notin\mathfrak a$.  Choose finitely many generators of
this kernel and multiply the associated $s$'s.
The resulting $y$ has $x(y)\ne0$ and annihilates the entire
kernel, exactly as required.
If a different original ring and lift representing the same
condition are specified, their two universal maps are inverse:
each composite fixes every representation entry and hence
is the identity by uniqueness.  Transporting the constructed
$x,\phi,\ker\phi,y$ through this unique isomorphism gives
the same completion and matrix-entry statements on that
specified ring.  Thus the chart conclusions hold for every
universal ring in \texttt{spChartL} and \texttt{spChartP},
not just for the presentation used in the construction.
\end{proof}
% END DX5

% BEGIN DX6
\begin{proposition}[DX6: the ordinary fourteen-parameter chart and \texttt{tanLz}]
\label{claim:DX6}
For the original good datum, the integral universal
$\operatorname{eqL}$ ring has the chart in DX5 with counts
$n(j)=d_L(j)$.  Thus its characteristic-zero germ is a formal
power-series ring in fourteen parameters with the exact weights
in (DX.4).

The following two assertions are also literal consequences
on the original $O$-valued tangent sets:
\[
\begin{aligned}
 z\in\operatorname{tanSet}_\ell(\operatorname{eqL},-p)
 &\ \Longrightarrow\
 \exists\,0\ne a_0\in O,\ m_0\in O:\\
 &\qquad a_0z(g)=(\Psi(g)^{-1}m_0-m_0)X_-
                        \quad(g\in G_\ell),\\
 z\in\operatorname{tanSet}_\ell(\operatorname{eqL},0)
 &\ \Longrightarrow\ z(s)=0\quad(s\in I_\ell).
\end{aligned}
 \tag{DX.29}
\]
Here $X_-=E_{20}+E_{31}$.
This is the complete \texttt{tanLz} statement, including its
single integral denominator and all original local elements.
\end{proposition}

\begin{proof}
Let $B=E\oplus\mathfrak n$ be a finite local augmented
$E$-algebra and let $\rho$ reduce to $r_0$.
Put $\lambda_i=\operatorname{ps}_i(s_0)\in O$.
AX6 makes these pairwise distinct in $E$.
Consequently their nonzero differences are units in $B$.
The characteristic-polynomial condition in (DX.25) yields
orthogonal idempotents
\[
 e_i=
 \frac{\prod_{h\ne i}(\rho(s_0)-\lambda_h)}
                 {\prod_{h\ne i}(\lambda_i-\lambda_h)}.
 \tag{DX.30}
\]
They sum to one and their images are free lines, since
their reductions are the four coordinate projections.
The stability equations say that
$F_j=\bigoplus_{i\le j}\operatorname{im}e_i$ is an invariant
complete flag.  The inertia equations say that $I_\ell$
acts on its $i$th successive quotient by exactly
$\operatorname{ps}_i$.

This flag is isotropic and self-dual for $J$.
Indeed preservation of the form at $s_0$ implies
$(\lambda_i\lambda_j-\theta(s_0))\langle u_i,u_j\rangle=0$
for vectors in its eigenlines.
For $i+j\ne3$ the scalar is nonzero, and hence a unit in $B$:
the four exponents $e+1,e,1,0$ have sum $e+1$ precisely
in the two complementary pairs, since $e\ge2$.
This proves the asserted orthogonality and flag duality.

The flag reduces to the standard one.
There is a unique element $u$ of the opposite unipotent
subgroup of $\operatorname{Sp}_4(B)$, reducing to $1$,
which takes the standard flag to it.
This is the ordinary Gaussian-elimination chart on isotropic
flags: choose the unit pivots in successive columns and use
the pairing to determine the complementary columns.
Equivalently, multiply the four lower-root subgroups from
(DX.3) in a fixed root order; solving the four lower
coordinates successively gives the unique element.
It has one parameter of each weight $-a,-b,-p,-t$.
These operations use only unit pivots, so apply over every
such $B$ and commute with its quotients.
The representation $u^{-1}\rho u$ is upper triangular.

Conversely, take any near-identity point $u$ of that
opposite-unipotent chart and any upper triangular
symplectic-similitude representation $\tau$ with multiplier
$\theta$ whose diagonal characters restrict to
$\operatorname{ps}_i$ on $I_\ell$.
Then $\rho=u\tau u^{-1}$ satisfies (DX.25).
For an upper triangular matrix with distinct diagonal
$\lambda_i$, its ordered spectral flag is the standard flag;
the diagonal action on each successive quotient is exactly
the required inertia character.  Applying (DX.30) therefore
recovers precisely the chosen flag and its unique $u$.
This proves an equivalence of the actual deformation
functors with a product of the four-parameter flag chart
and the upper triangular deformation functor.
It does not choose an arbitrary full eigenbasis or introduce
extra upper triangular changes of basis.

We prove the upper triangular functor unobstructed.
Its four diagonal characters have the form
$\operatorname{ps}_i\delta_i$, where $\delta_i$ is unramified
and congruent to $1$ modulo $\mathfrak n$, and
\[
                 \delta_0\delta_3=\delta_1\delta_2=1.
 \tag{DX.31}
\]
There are exactly two independent such characters.
On $G_\ell/I_\ell=\widehat{\mathbb Z}$ each is determined
by an arbitrary Frobenius value in $1+\mathfrak n$.
The logarithm and exponential are finite polynomials on
this nilpotent ideal; the associated additive homomorphism
factors through the natural projection
$\widehat{\mathbb Z}\to\mathbb Z_\ell$.
This proves existence, continuity, and the lifting of each
diagonal character across an arbitrary Artinian quotient.

For each positive root character
\[
                    \kappa\in\{\alpha,\beta,\Psi,\theta\}
 \tag{DX.32}
\]
we have $\kappa\ne1$ and $\kappa\ne\chi$ on $G_\ell$.
Nontriviality follows at $s_0$: their respective exponents
are $1,e-1,e,e+1$, all positive.
For $\alpha,\beta,\theta$, residual separation at the
actual $\sigma_\ell$ gives
$\kappa(\sigma_\ell)\ne1=\chi(\sigma_\ell)$.
For $\Psi$, comparison at $s_0$ gives $\Psi\ne\chi$
because $e\ge2$ and $\chi(s_0)$ has infinite order.
Thus DX1 gives
\[
 H^0(G_\ell,E(\kappa))=H^2(G_\ell,E(\kappa))=0,\qquad
 \dim_EH^1(G_\ell,E(\kappa))=1.
 \tag{DX.33}
\]
In particular the continuous cocycle space has dimension
two: its coboundaries form a one-dimensional copy of
$E(\kappa)$.

Consider a small extension $B\to B'$ with kernel $I$,
where $\mathfrak n_B I=0$.
Lift the two unramified diagonal characters first.
The upper unipotent group has four polynomial root
coordinates and is an affine four-space.
Choose continuous $E$-linear lifts of the coordinate
functions of a given representation over $B'$.
Together with the lifted diagonal they form a continuous
matrix-valued lift, preserving the multiplier and the
inertia diagonals, though initially it need not multiply
correctly.
Its multiplicative defect is a continuous two-cocycle with
values in
\[
 I\otimes_E\mathfrak n^+
   =\bigoplus_{\kappa\ {\rm in}\ (DX.32)}
                              I\otimes_EE(\kappa).
 \tag{DX.34}
\]
Associativity gives the cocycle equation.  Since $I$ is
killed by the maximal ideal, the conjugation action here
is exactly that of $r_0$, not a varying deformed action.
Equation (DX.33) makes the defect a continuous coboundary.
Correct the lifted matrices by the corresponding
$1$-cochain in $I\otimes\mathfrak n^+$; the square-zero
multiplication formula then removes the defect.
The correction is upper unipotent, so preserves the
multiplier, diagonals and inertia characters.
This proves lifting across every small extension.
Factoring the kernel of any Artinian surjection by its
finite maximal-ideal filtration proves lifting for every
surjection.

The tangent computation is now explicit.
The diagonal directions are the two unramified characters,
both weight zero.
The four positive-root directions are the cocycles in
(DX.33), two of each of the weights $a,b,p,t$.
There is no quotient by coboundaries: these are actual
matrix representations.  Finally flag conjugation supplies
one direction at each negative root.
The total is $2+4\cdot2+4=14$, and the weight multiplicities
are exactly $d_L$.
The entry action of the diagonal torus preserves the
flag chart, the Borel subgroup and all the constructions;
its tangent characters are precisely the entry weights
in (DX.3).  DX5 therefore supplies the claimed weighted
formal chart, with the original origin $r_0$.

To obtain (DX.29), start with an actual $O$-valued tangent
$z$ and its dual-number lift from \texttt{tanSet}, and
extend coefficients injectively to $E[\epsilon]$.
At weight $-p$ the preceding product decomposition has
only a flag direction.  Infinitesimal conjugation by
$1+\epsilon mX_-$ changes $r_0(g)$ by the cocycle
$(1-\Psi(g)^{-1})mX_-$.
Changing the sign of $m$ yields
$z(g)=(\Psi(g)^{-1}-1)m_E X_-$ for one $m_E\in E$,
simultaneously for all $g$.
Choose $0\ne a_0\in O$ with $m_0=a_0m_E\in O$.
Injectivity of $O\hookrightarrow E$ proves the first
identity on the original $O$-valued matrices.
At weight zero the only directions are the two unramified
diagonal characters, so the tangent vanishes on every
element of the actual inertia subgroup.
Injectivity again returns that equality to $O$.
These are exactly the two quantified conclusions in
\texttt{SpWApply.tanLz}.
\end{proof}
% END DX6

% BEGIN DX7
\begin{proposition}[DX7: the original finite-inertia ten-parameter chart]
\label{claim:DX7}
For each original $q\in S_b$, the integral universal
$\operatorname{eqP}(q)$ ring has the chart in DX5 with counts
$n(j)=d_P(j)$.  Its characteristic-zero germ has ten parameters,
two of weight zero and one of each root weight.
The statement imposes the source's entire inertia-kernel
condition on every deformation.
\end{proposition}

\begin{proof}
Set $N=I_q\cap\ker\psi$.  This is closed and normal in
$G_q$, since $\psi$ is a character and inertia is normal.
Every $\operatorname{eqP}$ deformation is trivial on $N$,
by its defining equations (DX.24), even when the coefficient
ring has nilpotents.  Thus it factors through the actual
profinite quotient
\[
  1\longrightarrow C=I_q/N
    \longrightarrow G_q/N
    \longrightarrow\widehat{\mathbb Z}\longrightarrow1 .
 \tag{DX.35}
\]
The map induced by $\psi$ identifies $C$ with the finite
cyclic subgroup $\psi(I_q)\subset O^\times$.
Its finiteness follows from the open kernel in
\texttt{Good}; a finite subgroup of a field's multiplicative
group is cyclic.
It is nontrivial, since $\ell$ divides its order by the
definition of $S_b$.
Conjugation on $C$ is trivial: $\psi$ is abelian and is
injective on this quotient.
Choose a lift $\gamma$ of $1\in\widehat{\mathbb Z}$.
The continuous power map from $\widehat{\mathbb Z}$ to the
procyclic closure of $\langle\gamma\rangle$ composes to the
identity on $\widehat{\mathbb Z}$.
It therefore splits (DX.35), giving
$G_q/N=C\times\widehat{\mathbb Z}$.
This quotient argument uses the original $I_q$ and $N$,
not merely their residual images.

On inertia, the original cyclotomic and $\omega$ characters
are trivial and $\xi$ is unramified.  Hence
\[
               r_0|_C=\operatorname{diag}(\psi,\psi,1,1).
 \tag{DX.36}
\]
Let $B=E\oplus\mathfrak n$ be an augmented finite local
$E$-algebra.  Every deformation of (DX.36) with the fixed
multiplier is conjugate by a symplectic matrix congruent
to $1$ to (DX.36) itself.
Here averaging is legitimate over $E$, including when
$\ell$ divides $|C|$; it is not used over the integral
residue field.
Explicitly, set
\[
                 u_0=|C|^{-1}\sum_{c\in C}
                                      \rho(c)r_0(c)^{-1}.
 \tag{DX.37}
\]
This reduces to $1$, is invertible, and satisfies
$\rho(c)u_0=u_0r_0(c)$.
Put $J'=u_0^tJu_0$ and $H=J^{-1}J'$.
The two forms $J,J'$ transform under $r_0(C)$ by the
same multiplier $\psi$, so $H$ commutes with $r_0(C)$.
It is self-adjoint for $J$ and congruent to $1$.
The finite binomial series $H^{-1/2}$ is therefore
defined in the nilpotent ideal and commutes with $r_0(C)$.
Then
$u=u_0H^{-1/2}$ still intertwines, and
$u^tJu=J$.
This proves the claimed symplectic rigidity of the
finite inertia representation, with explicit formulas
compatible with all coefficient quotients.

The centralizer $H_C$ of (DX.36) in $\operatorname{Sp}_4$
is $\operatorname{GL}_2$.
To check this over every $B$, choose $c\in C$ with
$\psi(c)\ne1$.  This difference is a unit in $E$ and
hence in $B$.  A commuting matrix is block diagonal on
the two eigenspaces $\langle e_0,e_1\rangle$ and
$\langle e_2,e_3\rangle$.
Preservation of $J$ makes its second block the paired
inverse transpose of the first, which is otherwise an
arbitrary invertible $2$-by-$2$ matrix.
Its Lie algebra has the two zero-weight directions
and the directions $a,-a$ in (DX.3).

There is a six-parameter formal slice for
$\operatorname{Sp}_4/H_C$ at the identity, of weights
$\pm p,\pm b,\pm t$.
One can construct it without a smooth-quotient theorem.
Near the identity the Cayley transform
\[
               Y\longmapsto(1+Y/2)(1-Y/2)^{-1}
 \tag{DX.38}
\]
identifies the formal Lie algebra with the formal
symplectic group; the inverse is
$g\mapsto2(g-1)(g+1)^{-1}$.
Take the six indicated root spaces as a torus-stable
linear complement to $\operatorname{Lie}H_C$.
The product of their Cayley slice with $H_C$ has derivative
the identity on the direct sum of these two Lie spaces.
Solving successively in each power of a nilpotent
augmentation ideal gives a unique formal factorization.
Thus every conjugating $u$ in (DX.37) has a unique
representative in this six-dimensional slice modulo
$H_C$.
The construction is torus equivariant, since (DX.38)
and the direct-sum decomposition are.

After this unique slice conjugation, inertia is (DX.36).
Because $\gamma$ commutes with $C$, its matrix must
centralize (DX.36).
It has the prescribed multiplier $\theta(\gamma)$
and reduces to $r_0(\gamma)$.
It is therefore uniquely an element of
$r_0(\gamma)\widehat H_C$, where $\widehat H_C$
is the formal group of $H_C$ at one.
Conversely every such matrix defines a continuous
representation of $\widehat{\mathbb Z}$.
To verify continuity, its entries and inverse entries
belong to a finite local $O$-subalgebra of $B$:
adjoin the finitely many nilpotent entry differences to
$O$, using $r_0(\gamma),r_0(\gamma)^{-1}\in\operatorname{GL}_4(O)$.
Products of sufficiently many nilpotent generators vanish,
so this subalgebra is finite over $O$.
Its matrix unit group is compact and profinite.
The integer power map of the chosen matrix therefore
extends to $\widehat{\mathbb Z}$.
Its multiplier agrees with the original $\theta$
on the dense integer subgroup, hence everywhere.
Combining this homomorphism with (DX.36) gives the
desired continuous representation of the original quotient
(DX.35), and then of $G_q$.

We have thus constructed an actual equivalence of functors
\[
 \{\operatorname{eqP}\text{-lifts near }r_0\}
   \ \simeq\
 \{\text{six slice parameters}\}\times\widehat H_C .
 \tag{DX.39}
\]
Both factors lift across all finite Artinian surjections.
The four centralizer parameters have weights $0,0,a,-a$;
the six slice parameters have the remaining root weights.
Their total multiplicities are exactly $d_P$.
The equivariance and exact functor identification now
apply DX5, giving the asserted chart and all its
completion properties on the original integral ring.
\end{proof}
% END DX7

% BEGIN DX8
\begin{theorem}[DX8: the complete original \texttt{LocBase} installation]
\label{claim:DX8}
For every original good datum $D$, there exists an actual
$\operatorname{LocBase}D$ with all its original coefficient,
decomposition and inertia data.
Together with DX2 and DX4, this installs precisely
\[
 \Phi_\ell(D.Gam),\qquad D.\operatorname{Nums},\qquad
                  \operatorname{Nonempty}(\operatorname{LocBase}D).
 \tag{DX.40}
\]
The local base has the exact smoothness multiplicities
\[
                       d_L(j)+|S_b|d_P(j)
                           \qquad(j\in\mathbb Z^2).
 \tag{DX.41}
\]
Thus these three actual arithmetic inputs in the literal
$\operatorname{exModel}$ constructor require no remaining
bespoke supplier once the original good datum is supplied.
\end{theorem}

\begin{proof}
For $v=\ell$ use DX5--DX6's integral universal local ring
$R_\ell$ with its actual $\operatorname{eqL}$ representation.
For each original $q\in S_b$ use DX5--DX7's $R_q$
with its $\operatorname{eqP}(q)$ representation.
The prime $\ell$ does not belong to $S_b$ by definition.
All these rings are over the same original $O$, and their
origins give exactly the corresponding restriction of
the same $r_0$.

Take their completed coproduct over $O$.
Explicitly, write their finite weighted presentations
$R_v=O[[z_{v,i}]]/J_v$, use the disjoint union of all
variable sets, and define
\[
       B=O[[z_{v,i}:v\in T]]/
                     \Bigl(\sum_{v\in T}J_v\Bigr).
 \tag{DX.42}
\]
There are finitely many variables and factors.
The ideal in parentheses is proper, because evaluation
of all variables at zero gives a map to $O$.
It is closed and homogeneous.  Thus $B\in\mathcal C_O$,
and its union presentation $W$ has all variable weights
in the required set and has an origin.
For any target $A\in\mathcal C_O$, maps $B\to A$
are exactly tuples of maps $R_v\to A$:
the variable images lie in $\mathfrak m_A$, so all
power-series substitutions converge, and the defining
ideals give exactly the required relations.
The integral universal properties of the $R_v$ identify
these with the arbitrary original local representation
tuples satisfying \texttt{LocEqs}.
Both existence and uniqueness are for actual tuples,
without quotienting by conjugacy.
This proves the \texttt{univ} field of \texttt{LocBase}.

The images of the universal local representations give
its \texttt{rho} and \texttt{eqs} fields.
Their entries remain in the actual weight parts of $W$,
giving \texttt{hom}.
Its origin evaluates them at $r_0$, giving \texttt{at0}.
These assertions hold on every element of every original
$G_v$; no isomorphic substitute field or representation
is chosen.

Complete (DX.42) at its characteristic-zero origin.
The resulting ring is the completed tensor product over
$E$ of the corresponding completions of all $R_v$.
This follows directly from the disjoint-variable
presentation, or by reducing modulo successive powers
of its augmentation ideal.
By DX6--DX7 it is therefore
\[
                 E[[t_1,\ldots,t_{14+10|S_b|}]],
 \tag{DX.43}
\]
with the disjoint union of their parameter weights.
The multiplicity of each $j$ is (DX.41).
The final graded-subalgebra argument of DX5 applies to
this union presentation verbatim: constants agree with
the origin, each parameter has a separately detected
linear coefficient, and a zero tangent is in the
$O$-saturation of the actual square of the origin ideal
inside $W.\!A$.
Consequently (DX.43) supplies exactly
\[
 \operatorname{SmoothPt}
      \bigl(W.\mathrm{gr},\,W.\mathrm{origin}|_{W.A},
                          j\mapsto d_L(j)+|S_b|d_P(j)\bigr).
\]
This is the remaining \texttt{smooth} field.
All fields of the original record have now been constructed.
DX2 and DX4 provide the other two statements in (DX.40).
\end{proof}
% END DX8

\subsection{Scope and named classical leaves}

The full \texttt{CohomologyHyp} record packages more than is
needed by this installation.  DX1 proves its native local
finiteness, local duality and local Euler clauses, global
finiteness and global Euler clauses, and the projection of
LD3 to its normalized Poitou--Tate clause.
LK2--LK4 supply the corresponding finite native-$K$ clauses,
LD2 supplies its unramified orthogonality clause, and DX3
proves its full odd-prime Greenberg--Wiles clause.
The entire \texttt{BlochKatoLocal} field, including its
arbitrary local period-module models, is not asserted here.
Accordingly this section does not promote the full
\texttt{LocalPublishedTheorems} or
\texttt{CohomologyHyp} records.
Instead DX5--DX8 prove the local universal rings and
the exact required ordinary and finite-inertia charts
directly.  Equations (DX.40)--(DX.41) are unconditional
relative to the original \texttt{Good} datum and the
following named classical inputs.

\begin{itemize}
\item \textbf{DX-LT and DX-LE.}
Finite local Tate duality and the local Euler--Poincare
formula, with the Brauer invariant normalization:
J.\ S.\ Milne, \emph{Arithmetic Duality Theorems},
second edition (2006), Chapter I, Theorem 2.1,
Corollary 2.3 and Theorem 2.8, printed pages 26--31.
The source-to-classical match for arbitrary finite,
possibly nonfree coefficients is the one already made
in LK1--LK4; DX1 supplies the original decomposition-group
transport and the rational inverse-limit argument.
\url{https://jmilne.org/math/Books/ADTnot.pdf}.
\item \textbf{DX-GF and DX-GE.}
Global finiteness for finite modules whose order is
invertible on the ring of $S$-integers, and the global
Euler--Poincare formula:
Milne, the same book, Chapter I, Corollary 4.15,
printed page 62, and Theorem 5.1, printed page 67.
The real-place factor and the range of high-degree
vanishing are retained exactly in DX1.
The finite Poitou--Tate input is the already proved LD3,
with its explicit Chapter I, Theorem 4.10 specialization
and compatible invariants.
\item \textbf{DX-CM.}
The Noetherian completion theorem, including exactness
for finite modules and agreement modulo powers of the
ideal; and Krull intersection.
The Stacks Project, Algebra, Sections 10.96--10.97
and Lemma 10.51.4.
Only the completion/quotient assertions and the injectivity
of a Noetherian local ring into its completion are used
in DX5.  The deformation representation, its
unobstructedness and the weighted chart are proved above,
not retained as a separate deformation-theoretic leaf.
\url{https://stacks.math.columbia.edu/tag/00M9};
\url{https://stacks.math.columbia.edu/tag/0BNH};
\url{https://stacks.math.columbia.edu/tag/00IP}.
\end{itemize}

The original decomposition-to-completion equivalence and
its roots-of-unity compatibility are the same precisely
matched classical input already recorded in KB6 and LK1.
All original arithmetic character, inertia and residual
separation facts are AX6.  The incoming good datum itself
is constructed in the separate arithmetic-data branch.
No automorphy theorem, rigidity predicate, prime-counting
supplier or final RH/GRH conclusion is claimed in this
local installation.
